% =============================================================================
% Chapter 8 --- Limitations and Future Work
% Grounded in Context 5 §3 (limitations / technical debt) and §4 (roadmap).
% =============================================================================

\chapter{Limitations and Future Work}
\label{ch:limitations}

This chapter catalogs the limitations of the present work and proposes
a concrete engineering roadmap for follow-on verification effort. Every
limitation below is traceable to the actual \sv{.sv} sources or to
comments in the Context documents; no extra-repository device assumptions
are used.

\section{Limitations and Technical Debt}
\label{sec:lim-limitations}

\subsection{No Independent ANN Behavioral or Numerical Model}
\label{sec:lim-nomodel}

The testbenches implement a sparse 4D array
\sv{ann\_weight\_matrix[...]} updated when programming is detected, and
return \signame{weight\_read\_data} as a function of the
\signame{ann\_core\_word} decode. That is a
\emph{register-transfer-level mock}, not a memristor or MAC-array model.
The following are therefore untested by this abstraction:

\begin{itemize}
    \item Real timing of \signame{op\_done} relative to array
          settling --- the TBs use deterministic counters or immediate
          assertions for many paths.
    \item End-to-end numerical accuracy of an inference --- no
          reference MAC against golden vectors is implemented in RTL.
    \item Cross-talk, IR drop, and variability beyond the discrete
          4-bit readback injection patterns used for recovery
          coverage.
\end{itemize}

The scope of what is validated is explicitly \emph{control sequencing}
and \emph{digital packet semantics} to a first-order memristor
abstraction, not full accelerator accuracy~\cite{Ielmini2018RRAM,
Sebastian2020IMC}.

\subsection{\texttt{op\_done} is a Testbench-Controlled Artifact in Most Benches}
\label{sec:lim-opdone}

In \sv{controller\_host\_read\_reorder\_tb.sv} and the program/verify
TBs, \signame{op\_done} is generated with cycle-count heuristics when
certain \signame{pulses} and DUT states are active. The risk is that
if a silicon \signame{op\_done} is late, early, or pulsed differently,
PROG/ERASE sub-FSM timing in hardware may diverge from what the
SystemVerilog mock proved. The repository does not include min/max
delay characterization --- only the FSM's dependency on
\signame{op\_done} in RTL is real. Concrete proposals for assertions
that could guard against this are listed in \Cref{sec:future-sva}.

\subsection{\stateval{S\_RESET} is Unreachable from \stateval{S\_IDLE}}
\label{sec:lim-sreset}

\stateval{S\_IDLE} decodes \sv{CMD\_*} into \stateval{S\_PROGRAM},
\stateval{S\_ERASE}, \stateval{S\_READ}, or \stateval{S\_COLLECT\_DATA};
it never sets \signame{next\_state} to \stateval{S\_RESET}. The state
\stateval{S\_RESET} still exists, drives \signame{ann\_reset} for a
cycle, and unconditionally forces \stateval{S\_PROGRAM} on exit. The
sequential block retains legacy initialization of programming-related
counters on the condition
\signame{state}~=~\stateval{S\_IDLE}~\sv{\&\&}~\signame{next\_state}~=~\stateval{S\_RESET}
(for example, \signame{weight\_count\_reg}, \signame{matrix\_rows},
\signame{matrix\_cols}, \signame{buffer\_idx\_reg}), but because no
path sets that transition, this code is dead in the current graph.

\paragraph{Debt.}
Either remove the dead transition after an audit, or add a host-visible
reset command that \emph{does} enter \stateval{S\_RESET} if product
needs global buffer/core reset from idle.

\subsection{\signame{weight\_read\_en} is Assigned but Unused}
\label{sec:lim-weightreaden}

A local \signame{weight\_read\_en} is toggled in \stateval{S\_VERIFY}
(sub-state \sv{VERIFY\_IDLE}) but the signal is not connected to a
port and does not feed any other logic in the excerpted design. The
actual verify path samples \signame{weight\_read\_data} in
\sv{VERIFY\_CHECK}. \textbf{Debt:} delete the dead signal or connect
it to a future explicit enable for the analog readback path.

\subsection{Module Parameters \sv{ADDR\_WIDTH} / \sv{WEIGHT\_WIDTH} Not Driving Logic}
\label{sec:lim-params}

As noted in \Cref{sec:analysis-params}, these parameters are not used
in the \sv{ann\_controller} body. \textbf{Debt:} either connect them to
port widths or remove them from the public API to avoid misleading
integrators who assume a width change re-parameterizes the core.

\subsection{Direct-Address versus Legacy Buffer-Index Flow}
\label{sec:lim-dual-flow}

The RTL and its comments indicate two conceptual programming histories:
a \emph{direct-address} path that uses \signame{address\_reg}[5:0] for
buffer read/write in PROG and VERIFY, and \emph{legacy}
\signame{buffer\_idx\_reg} / \signame{weight\_count\_reg} support that
is ``kept for ERASE/VERIFY compatibility''. The signal
\signame{weight\_prog\_done} is asserted in both the legacy
``advance buffer on verify re-entry to program'' path and the direct
\sv{PROG\_COMPLETE} path. \textbf{Debt:} a maintenance pass could unify
or formally deprecate the unused path to reduce the double semantics.

\subsection{\signame{retry\_cnt} in the Direct-Address Mode}
\label{sec:lim-retry}

\sv{controller\_prog\_verify\_lut\_10w\_tb.sv} documents that the
\emph{erase max-retry give-up} path (\signame{erase\_max\_retries\_exhausted}
/ \signame{retry\_cnt} accumulation) is not reachable in the
direct-address flow as wired, because \sv{PROG\_COMPLETE} clears
\signame{retry\_cnt} when \signame{buffer\_idx\_reg}~\(<\)~\signame{weight\_count\_reg}~\(-\)~1
with the default 640 count. Erase retries therefore do not accumulate
across PROG cycles in that configuration. The limitation is specific
to the interaction of reset conditions and the default
\signame{weight\_count\_reg}; it is not a claim about all possible
parameterizations. A future test should either shrink
\signame{weight\_count\_reg} in a contrived scenario or drive the
legacy buffer sweep to hit \signame{erase\_max\_retries\_exhausted}.

\subsection{Repository Scope}
\label{sec:lim-scope}

There is no UVM testbench, no checked-in formal property set, and no
gate- or post-layout netlist in the explored tree. These absences are
out of scope for the present project, not claims of impossibility.
Suggested next steps are given in \Cref{sec:future-sva,sec:future-cr,sec:future-next}.

\section{Future Work --- SystemVerilog Assertions}
\label{sec:future-sva}

Assertions should bind to \modname{ann\_controller} instances in
testbenches (or a thin wrapper module) so that the RTL files remain
lint-clean. Using \sv{bind} is appropriate if the compile flow
supports it. Natural interfaces to watch are \signame{valid},
\signame{cmd}, \signame{busy}, \signame{pulses}, \signame{op\_done},
\signame{state}, \signame{prog\_state}, \signame{verify\_state}, and
\signame{erase\_state} (the DUT is already referenced hierarchically in
TBs, for example \signame{dut.state}).

\Cref{tab:future-sva} lists concrete checks grounded in
\Cref{ch:rtl}'s description of the RTL behavior.

\begin{table}[htbp]
    \centering
    \caption{Concrete SVA proposals, each grounded in a specific section
        of \Cref{ch:rtl}.}
    \label{tab:future-sva}
    \begin{tabularx}{\linewidth}{X l}
        \toprule
        \textbf{Concern and suggested check} & \textbf{Grounding} \\
        \midrule
        \textbf{Handshake.} If \signame{valid} rises while
            \signame{state}~=~\stateval{S\_IDLE} and a legal \signame{cmd}
            is sampled, \signame{busy} or the main state must leave
            \stateval{S\_IDLE} within a bounded number of cycles, unless
            \signame{rst\_n} drops. &
            \Cref{sec:rtl-controller-mainfsm} (\stateval{S\_IDLE} decode).
        \\ \midrule
        \textbf{\signame{op\_done} usage.} In \sv{PROG\_SELECT},
            \sv{PROG\_WAIT\_ACK}, \sv{ERASE\_SELECT}, and
            \sv{ERASE\_WAIT\_ACK}, simulation must not observe an \sv{X}
            on \signame{op\_done}; the assertion can be strengthened into
            a stability or toggle property once analog bounds are
            supplied by a device team. &
            \Cref{sec:rtl-controller-prog,sec:rtl-controller-erase}.
        \\ \midrule
        \textbf{Pulse / mode consistency.} When \signame{pulses} is a
            non-HIZ mode, \signame{pulse\_cnt} stays within
            \(0..\signame{pulse\_total}-1\) until \signame{pulse\_done}
            asserts (one property per active command state). &
            \Cref{sec:rtl-controller-pulses}.
        \\ \midrule
        \textbf{Invalid tail.} At the PI boundary,
            \sv{ann\_tail\_is\_valid\_onehot(host\_data[23:0])} is a
            precondition for \signame{valid}; mirror this as an assume
            in random-stimulus TBs and as a cover in directed ones. &
            \Cref{sec:rtl-pi-onehot}.
        \\ \midrule
        \textbf{INF handshake extension.} If INF is later extended to
            sample \signame{op\_done}, add cases where \signame{op\_done}
            is not tied low, unlike the present
            \sv{controller\_inf\_buffer\_flow\_tb}. &
            \Cref{sec:analysis-inf-opdone}.
        \\ \bottomrule
    \end{tabularx}
\end{table}

\paragraph{Integration point.}
The most natural first targets are the PI integration bench and the
program--verify LUT bench: they already instantiate the DUT and
reference internal state for \signame{op\_done} generation, so binding
an assertion module introduces no new plumbing.

\section{Future Work --- Constrained-Random Stimulus and Coverage}
\label{sec:future-cr}

The existing infrastructure already provides a \emph{legal-packet
generator}: \sv{build\_host\_ann\_word} and
\sv{ann\_tail\_is\_valid\_onehot} together can generate and filter
random host packets. A concrete next step is:

\begin{enumerate}
    \item \textbf{Randomize.} Randomize \signame{host\_data}[31:0] and
          \signame{host\_cmd} subject to one-hot tail validity and legal
          \signame{cmd} values (from \sv{cmd\_t} in
          \sv{parallel\_interface\_pkg}).
    \item \textbf{Constrain.} Constrain sequences to the controller
          FSM's protocol: do not send a new \signame{valid} while
          \signame{busy} is high unless the protocol is extended to
          allow pipelining. The current FSM does not, so \emph{idle-wait}
          is the safe constraint.
    \item \textbf{Cover.} Define covergroups on \signame{state},
          \signame{prog\_state}, \signame{verify\_state},
          \signame{erase\_state}, and on the three branches of
          \sv{VERIFY\_CHECK} (match, under-programmed, over-programmed),
          together with crosses of \signame{cmd} at idle.
    \item \textbf{Self-checking.} Reuse the shadow-matrix pattern from
          \sv{controller\_prog\_verify\_lut\_tb.sv} to compare
          \signame{ann\_core\_word} decode to expected storage, and
          keep the directed regressions as a shrinking baseline to
          which random-found failures can be added.
\end{enumerate}

\section{Future Work --- Recovery and Corner-Case Expansion}
\label{sec:future-recovery}

The intentional \sv{read<expected} and \sv{read>expected} injection
should be parameterized (indices, counts) so that
\signame{prog\_retry\_cnt}~=~\sv{MAX\_PROG\_RETRIES} and multiple ERASE
rounds can be reached without relying on a 640-weight runtime every
time. In addition, host-initiated erase and verify-initiated erase
should both be covered, observing the \signame{erase\_from\_host}
register in waveforms to confirm which path is taken.

\section{Future Work --- Consolidation and Documentation Alignment}
\label{sec:future-next}

\paragraph{Consolidation of programming modes.}
If the product direction settles on a single programming interface
(buffer-sweep only, or direct only), the dead \stateval{S\_RESET}
path, the unused module-level parameters, and the dual
\signame{weight\_prog\_done} semantics can be refactored against a
written spec. This reduces the fan-out of any future SVA or
constrained-random environment built on top.

\paragraph{Documentation alignment.}
The repository README uses the phrase \emph{in-memory computing} (IMC)
at the project level; RTL comments refer to the parallel interface and
the ANN core. The glossary in \Cref{app:fsm-diagrams} and the rest of
the report relate those terms without attributing IMC behavior to a
specific SystemVerilog module name, so that new readers are not misled
into looking for an \sv{imc} module that does not exist in-tree.

\paragraph{Gate-level flow.}
Finally, once the RTL and functional verification stabilize, the
natural next step in a VLSI flow is synthesis, static timing analysis,
and power/area characterization. Those steps require a target cell
library, clock/reset/timing constraints, and a top-level integrating
module, none of which are in scope of the present project.
